\section{Related Work}
\label{sec:related}

% The subsections position the shared fragment representation against its
% neighboring research lines. They are not a capability-scorecard.

% ── Versioning filesystems ─────────────────────────────────────────────────
\subsection{Versioning Filesystems}
\label{subsec:rw-versioning}

Prior versioning filesystems fall into three broad families. \emph{Per-write /
comprehensive} systems
keep a version on every write: Elephant retains all versions and lets
user policy decide what to forget~\citep{Santry1999Elephant}. CVFS versions at
write granularity with space-efficient multiversion metadata for a fixed
security window~\citep{Soules2003CVFS}. Wayback and VersionFS realize the same
idea as a user-level undo log and a stackable
layer~\citep{Cornell2004Wayback, MuniswamyReddy2004VersionFS}. \emph{Snapshot /
copy-on-write} systems version at point-in-time, filesystem-wide granularity:
WAFL duplicates the root inode~\citep{Hitz1994WAFL}, ZFS and Btrfs snapshot
whole datasets/subvolumes over CoW trees~\citep{Rodeh2013Btrfs}, and ext3cow
adds epoch snapshots to ext3~\citep{Peterson2005Ext3cow}. \emph{Content-addressed
/ git-like} systems version through immutable objects and discrete commits: Ori
gives a filesystem a git-style replicated, integrity-checked
history~\citep{Mashtizadeh2013Ori}, and ForkBase pulls git-conformant branching
into a storage substrate~\citep{Lin2020ForkBase}.

\sfs{} differs from all three, but does not claim to be the first filesystem
with native multiversion structures. In particular, CVFS uses journal-based
multiversion B-trees for filesystem metadata~\citep{Soules2003CVFS}. In \sfs{},
versioning is neither an overlay
(Wayback/VersionFS), a filesystem-wide snapshot (WAFL/ZFS/Btrfs/ext3cow), nor a
separate object store (Ori/ForkBase/git): each unit carries its own
continuously advanced fragment lineage as the \emph{native} on-disk model, with
named commits as retention-protected fixpoints layered over the running history
(\cref{subsec:mvcc}). Multiversion concurrency control is well established in
databases. Recent work extends it to disaggregated
memory~\citep{Zhang2024Motor}. The narrower gap \sfs{} targets is a per-file
fragment version map that is simultaneously the unit of history,
synchronization, and sub-file conflict handling, combined with the encrypted
multi-writer and cross-surface container design.

% ── CoW, crash consistency, allocator ──────────────────────────────────────
\subsection{Copy-on-Write, Crash Consistency, and Allocation}
\label{subsec:rw-cow}

Crash-consistency mechanisms split into journaling, explicit dependency
ordering, and atomic root publication.  Journaling ranges from
ext3/4 to optimizations that decouple ordering from durability
(OptFS~\citep{Chidambaram2013OptFS}). Ordering approaches include soft
updates~\citep{Ganger2000SoftUpdates}, Featherstitch's generalized
dependencies~\citep{Frost2007Featherstitch}, and, recently, SquirrelFS, which
encodes update ordering in Rust's typestate and checks it at compile
time~\citep{LeBlanc2024SquirrelFS}. They avoid a journal but have no single atomic
point. \sfs{} uses the third mechanism for publication: the committed root is a
sequence-wins choice between two alternating header slots. This is not a
journal-free-system claim: an optional WAL serves asynchronous writes, and
in-place overwrite recovery uses undo records in the eviction tail. WAFL's two
alternating \texttt{fsinfo} blocks~\citep{Hitz1994WAFL} are the closest
precursor to the root publication. ZFS's uberblock ring uses the same principle.
Our novelty is not that commit pattern but its combination with
\emph{fragment-granular} CoW over a flat three-region allocator: WAFL, ZFS, and
Btrfs shadow tree paths recursively up to the root,
which measurably amplifies writes~\citep{Rodeh2013Btrfs}. \sfs{} does not place
file content in such a filesystem-wide CoW tree, although its key and identity
catalog tries still incur their own CoW update cost
(\cref{subsec:commit,subsec:alloc}). On space management, LFS and
its flash/NVM successors (F2FS~\citep{Lee2015F2FS}, NOVA~\citep{Xu2016NOVA}) are
log-structured with a cleaner/GC.  \sfs{} has no background wandering cleaner:
fixed regions (catalog head, live units, eviction tail) replace the wandering
log, while retention expiry, reclamation, defragmentation, and trim are explicit
operations.  The eviction-tail region has no direct analogue in that
literature. WineFS's finding that alignment and aging affect fragmentation
substantially~\citep{Kadekodi2021WineFS} reinforces the need to evaluate the
long-term behavior of this simpler allocator separately.

% ── Chunking ────────────────────────────────────────────────────────────────
\subsection{Chunking: Content-Defined versus Fixed-Size}
\label{subsec:rw-chunking}

Content-defined chunking (CDC), introduced by LBFS~\citep{Muthitacharoen2001LBFS}
and since accelerated by Data Domain~\citep{Zhu2008DataDomain} and the FastCDC
line~\citep{Xia2016FastCDC}, sets chunk boundaries by a rolling hash so that an
inserted byte does not shift all following boundaries. The trade-off is
content-dependent boundaries and a rolling-hash scan. CDC is deterministic for
a fixed input and algorithm. Its boundaries are content-dependent, not random.
\sfs{} deliberately rejects CDC in favor of fixed-size fragments with a
$\sqrt{\cdot}$-shaped schedule: within a size band, fixed boundaries give a
constant-time position$\rightarrow$fragment mapping (fragment $n$ of a version
maps to fragment $n$ of the next), with no rolling-hash scan. A sync round still
scans the unit manifests and fragment-version vectors, so its metadata work is
linear in the number of units and fragments. Normal steady-state sync transfers
changed or missing fragments. At a fragment-size band transition, \sfs{}
re-chunks the unit.
Advances in CDC \emph{throughput} do not change this content-dependence trade-off.
The Merkle-/content-addressed tradition (Venti~\citep{Quinlan2002Venti},
git, IPFS~\citep{Benet2014IPFS}, Tahoe-LAFS~\citep{WilcoxOHearn2008Tahoe}) informs
\sfs{}'s Merkle-like catalog trie, but \sfs{} keeps identity+version addressing
for data blocks rather than pure content addressing, precisely because it
synchronizes mutable versions rather than stacking immutable snapshots.

% ── Replication and conflict resolution ────────────────────────────────────
\subsection{Optimistic Replication and Conflict Resolution}
\label{subsec:rw-replication}

\sfs{}'s synchronization descends from version vectors~\citep{Parker1983VersionVectors}
and the optimistic-replication filesystems built on them: Coda's disconnected
operation~\citep{Kistler1992Coda}, Ficus's peer replication with VV conflict
detection~\citep{Reiher1994Ficus}, Bayou's application-defined merge
procedures~\citep{Terry1995Bayou}, and Ivy's per-participant logs in a
DHT~\citep{Muthitacharoen2002Ivy}. Across this family, \emph{content} conflicts
within a file are left whole-file and usually manual (Ficus auto-merges only
directory structure, and Ivy needs a conflict tool). At the other pole, CRDTs
eliminate conflicts algebraically~\citep{Shapiro2011CRDT}: JSON-CRDT/Automerge
converge nested documents~\citep{Kleppmann2017JSONCRDT}, ElmerFS models
filesystem structures as CRDTs over a CRDT
store~\citep{Elmerfs2021}, at the price of per-element metadata,
tombstones, and sometimes counter-intuitive auto-resolutions. \sfs{} sits
between the two: version-vector sync with \emph{block-granular} auto-merge.
With strain splits, concurrent edits to different fragments of the same file
merge automatically, while concurrent updates to the same fragment remain a
conflict even when the changed byte ranges inside it do not overlap.  The
conflict is surfaced as a \texttt{user.sfs.conflict} xattr. This fills the gap
the VV filesystems leave (sub-file content merge) without adopting CRDT
machinery.

% ── Portable encrypted VFS and cloud substrates ────────────────────────────
\subsection{Portable Encrypted VFS and Cloud Substrates}
\label{subsec:rw-vfs}

WNFS is the closest architectural neighbor on the web-facing side: it is a
versioned, content-addressed, encrypted filesystem designed for offline
collaboration, untrusted storage, and WebAssembly, and models the tree as a
state-based CRDT~\citep{WNFSspec}.  Peergos likewise combines client-side
encryption, protected metadata, capability sharing, cloud or peer storage, and
web and desktop clients~\citep{Peergos}.  These systems demonstrate that an
encrypted application-owned filesystem is a useful web abstraction.  \sfs{}
does not claim that category.  Its narrower distinction is to carry one
position-addressed fragment/version representation through ordinary file APIs,
a raw-device kernel implementation, FUSE, an application VFS, and an
experimental wasm binding.  It preserves concurrent same-fragment branches
rather than making the entire tree a CRDT.

At the infrastructure scale, OceanStore treats persistent encrypted objects as
the basis for global storage~\citep{Kubiatowicz2000OceanStore}, while Ceph
builds file, block, and object services above a distributed object
store~\citep{Weil2006Ceph}.  \sfs{} borrows the unifying-substrate ambition but
not their scale claim: its hosted service is an untrusted relay for the same
client-owned records used locally, and the present evaluation contains neither
a storage cluster nor a multi-client scalability experiment.  The result is a
backend-neutral VFS boundary, not a replacement for a distributed object-store
control plane.

% ── Cryptographic / untrusted storage ──────────────────────────────────────
\subsection{Cryptographic and Untrusted Storage}
\label{subsec:rw-crypto}

Cryptographic filesystems progress by threat model. Farsite is an important
early combination: it encrypts file data placed on incompletely trusted desktop
machines and uses Byzantine fault tolerance for distributed filesystem
metadata~\citep{Adya2002Farsite}. Local, honest-but-curious
(HbC) systems such as CFS~\citep{Blaze1993CFS}, \texttt{gocryptfs}/eCryptfs,
and dm-crypt/LUKS give at-rest confidentiality for one user, without multi-writer
or sync. Secure-sharing-on-untrusted-storage systems add key management and
revocation: Plutus with lazy revocation via key
rotation~\citep{Kallahalla2003Plutus}, SiRiUS with signed per-file metadata over
an untrusted underlay~\citep{Goh2003SiRiUS}, Cryptree with hierarchical key
derivation~\citep{Grolimund2006Cryptree}, and Key Regression formalizing the
rotation~\citep{Fu2006KeyRegression}. A stronger model assumes an actively
byzantine server and provides \emph{fork consistency}: SUNDR~\citep{Li2004SUNDR}
makes any server equivocation detectable, and SPORC extends this to collaborative
multi-writer apps~\citep{Feldman2010SPORC}. Depot tolerates faulty nodes and
heals forks~\citep{Mahajan2011Depot}. Recent metadata-hiding systems push
further with ORAM (Metal~\citep{Chen2020Metal}, Titanium~\citep{Chen2022Titanium})
or anonymity (Ghostor~\citep{Hu2020Ghostor}), and KBFS signs a recursive root
block for integrity under total server compromise. Cryptanalysis of deployed
end-to-end-encrypted storage~\citep{Backendal2023MEGA, Backendal2024BrokenEcosystem}
motivates the whole line: real providers turn malicious, and the server must
never see keys or paths.

\sfs{} takes the \emph{honest-but-curious} position deliberately, not the
fork-consistency one.  In its secure hosted profile, metadata and content are
client-side AES-GCM encrypted and accepted record frontiers are authorized by
Ed25519 signatures under an owner-signed, epoch-versioned writer set.  The
server receives neither keys, paths, metadata plaintext, nor content plaintext.
This statement does not extend to the selectable XTS or plaintext content
modes: XTS provides confidentiality without authentication, and \textsc{none}
provides neither.  We use \emph{client-side-encrypted storage} rather than the
cryptographic term zero knowledge because the service still observes accounts,
object identifiers, version vectors, counts and sizes, timing, and access
patterns.  The signatures cover logical record frontiers, not a complete
per-fragment provenance proof.

Multi-writer authorization comes from the signed writer set (tombstone
revocation and root-key rotation) rather than from a
fork-consistency protocol or CRDT/OT. Relative to Plutus, revocation is bound to
epochs and the writer set rather than to per-file-group key chains. Relative to
SUNDR/SPORC, \sfs{} does not guarantee detection of a server that withholds or
forks versions. This is an explicit model weakening we defend in \cref{sec:discussion}.
KBFS is a close deployed neighbor~\citep{KeybaseKBFS}. \sfs{} differs in
binding identity to SRP-6a password authentication and a self-contained
container rather than to social sigchains, and in targeting one self-contained
format across kernel, FUSE, and experimental wasm bindings.

Finally, the earlier Self-Certifying File System also used the acronym
SFS~\citep{Fu2000SFS}.  It secured remote access through self-certifying
pathnames and is otherwise unrelated to the artifact described here.  This
paper consistently typesets the present system as \sfs{}.
